\section{Loss Optimization Suite and Training Dynamics}
\label{sec:experiments}

\subsection{Comprehensive 7-Term Loss Suite}
To achieve high radiometric precision while avoiding the classic blurriness of L1 autoencoders, the generator is supervised under a composite 7-term loss suite:

\subsubsection{Heteroscedastic Gaussian NLL Loss ($\mathcal{L}_{\text{NLL}}$)}
Following Kendall and Gal \cite{kendall2017what}, we supervise the predicted mean $\hat{\mathbf{y}}$ and log-variance $\mathbf{s} = \log \sigma^2$ using an observation-noise attenuated formulation:
\begin{align}
    \mathcal{L}_{\text{NLL}} &= \frac{1}{2HWC} \sum_{i,j,c} \Big( \exp(-\mathbf{s}_{i,j,c}) \|\hat{\mathbf{y}}_{i,j,c} - \mathbf{y}_{i,j,c}\|_1 \nonumber \\
    &\qquad\qquad\quad + \mathbf{s}_{i,j,c} \Big)
\end{align}
The exponential term dynamically downweights loss gradients in intrinsically ambiguous, heavily occluded pixels while the linear $\mathbf{s}$ penalty prevents variance explosion ($\mathbf{s} \to \infty$).

\subsubsection{Multi-Scale Spectral L1 Loss ($\mathcal{L}_{\text{L1}}$)}
Direct absolute pixel difference enforces low-frequency radiometric fidelity across visible and near-infrared bands:
\begin{equation}
    \mathcal{L}_{\text{L1}} = \frac{1}{HWC} \|\hat{\mathbf{y}} - \mathbf{y}\|_1
\end{equation}

\subsubsection{Spectral-Normalized PatchGAN Adversarial Loss ($\mathcal{L}_{\text{Adv}}$)}
We employ a two-scale discriminator $D = \{D_1, D_2\}$ operating on full and half resolutions. All convolutional layers in $D$ are stabilized using Spectral Normalization \cite{miyato2018spectral} to enforce the 1-Lipschitz condition. Under the Least-Squares GAN (LSGAN) formulation:
\begin{equation}
    \mathcal{L}_{\text{Adv}} = \frac{1}{2} \sum_{k=1}^2 \mathbb{E}_{\mathbf{X}}\left[ (D_k(\hat{\mathbf{y}}) - 1)^2 \right]
\end{equation}

\subsubsection{Discriminator Feature Matching Loss ($\mathcal{L}_{\text{FM}}$)}
Feature matching minimizes the L1 distance between intermediate activation representations of real and synthesized images across all discriminator layers:
\begin{equation}
    \mathcal{L}_{\text{FM}} = \sum_{k=1}^2 \sum_{l=1}^{L_k} \frac{1}{N_{k,l}} \|D_k^{(l)}(\mathbf{y}) - D_k^{(l)}(\hat{\mathbf{y}})\|_1
\end{equation}

\subsubsection{Frequency-Domain FFT Magnitude Spectrum Loss ($\mathcal{L}_{\text{FFT}}$)}
To penalize spatial oversmoothing in the frequency domain, we compute the 2D Fast Fourier Transform (FFT) per spectral band:
\begin{align}
    \mathcal{L}_{\text{FFT}} &= \frac{1}{C} \sum_{c=1}^C \Big\| \log\left(1 + |\mathcal{F}(\hat{\mathbf{y}}_c)|\right) \nonumber \\
    &\quad - \log\left(1 + |\mathcal{F}(\mathbf{y}_c)|\right) \Big\|_1
\end{align}

\subsubsection{Multi-Scale Sobel Structural Edge Loss ($\mathcal{L}_{\text{Sobel}}$)}
To enforce sharp gradient consistency along structural boundaries, Sobel convolutions extract spatial edge magnitude maps $\mathbf{G}$:
\begin{equation}
    \mathcal{L}_{\text{Sobel}} = \frac{1}{C} \sum_{c=1}^C \|\mathbf{G}(\hat{\mathbf{y}}_c) - \mathbf{G}(\mathbf{y}_c)\|_1
\end{equation}
where $\mathbf{G}(\mathbf{y}) = \sqrt{(\mathbf{K}_x * \mathbf{y})^2 + (\mathbf{K}_y * \mathbf{y})^2}$ via horizontal and vertical kernels $\mathbf{K}_x, \mathbf{K}_y$.

\subsubsection{Spectral Angle \& Vegetation Index Consistency Loss ($\mathcal{L}_{\text{Spec}}$)}
To guarantee radiometric preservation for downstream agricultural analysis, we penalize Normalized Difference Vegetation Index (NDVI) deviation:
\begin{align}
    \text{NDVI}(\mathbf{y}) &= \frac{\mathbf{y}_{\text{NIR}} - \mathbf{y}_{\text{Red}}}{\mathbf{y}_{\text{NIR}} + \mathbf{y}_{\text{Red}} + \epsilon} \\
    \mathcal{L}_{\text{Spec}} &= \|\text{NDVI}(\hat{\mathbf{y}}) - \text{NDVI}(\mathbf{y})\|_1
\end{align}

\subsubsection{Total Generator Objective \& Loss Weighting Rationale}
The overall multi-task training objective is defined as:
\begin{align}
    \mathcal{L}_{\text{Total}} &= 1.0\,\mathcal{L}_{\text{NLL}} + 10.0\,\mathcal{L}_{\text{L1}} + 10.0\,\mathcal{L}_{\text{FM}} \nonumber \\
    &\quad + 0.1\,\mathcal{L}_{\text{FFT}} + 0.5\,\mathcal{L}_{\text{Sobel}} + 0.2\,\mathcal{L}_{\text{Spec}} + 1.0\,\mathcal{L}_{\text{Adv}}
\end{align}
The dominant weighting assigned to $\mathcal{L}_{\text{L1}}$ ($w_{\text{L1}} = 10.0$) and feature matching relative to the adversarial penalty ($w_{\text{adv}} = 1.0$) is deliberately designed to enforce strict radiometric preservation of physical surface reflectance (vital for agricultural NDVI analytics), while the FFT and Sobel edge losses eliminate texture oversmoothing without unconstrained GAN hallucination.

\subsection{Distributed Training Implementation Details}
The model was trained for 50 epochs on dual NVIDIA Tesla T4 GPUs (16\,GB VRAM each) using PyTorch Distributed Data Parallel with NCCL backend and Automatic Mixed Precision (\texttt{fp16}). Optimization utilized the AdamW optimizer ($\beta_1 = 0.5$, $\beta_2 = 0.999$, weight decay $10^{-4}$) with a batch size of 16 patches per GPU (effective batch size 32). The learning rate was initialized at $2 \times 10^{-4}$ and decayed to $10^{-6}$ via a cosine annealing schedule. An Exponential Moving Average (EMA) shadow model with decay $\beta_{\text{EMA}} = 0.999$ was maintained for all inference evaluations. Total training completed in 4.7 hours across 11,950 training patches.
